\section{Experimental protocol}
\label{sec:protocol}

\subsection{Recipe}

Every run in this paper trains the same pinned configuration, summarized in
Table~\ref{tab:recipe}. It is a 32-layer, 512-wide Llama-style transformer with
pre-RMSNorm blocks, separate $q,k,v,o$ projections and a SwiGLU MLP at
expansion 4, totalling 265.6M trainable parameters of which 134.3M are the
block matrices the optimizer modifies. Training consumes 2{,}684{,}354{,}560
tokens of a DCLM/StarCoder/Proof-Pile mixture in 5{,}120 updates at a global
batch of 524{,}288 tokens, on eight A100-SXM4-40GB GPUs with DDP. The schedule
holds the learning rate flat for 1{,}024 updates and then decays it linearly,
with no warmup.

Two aspects of the protocol do more work than usual and are worth stating
plainly. First, data order is \emph{pinned by seed}: global sequence
identifiers assign the same tokens to the same update index across all arms, so
two runs at one seed are exactly paired update by update. Second, the first 64
permuted windows of each training source are excluded from training and
reserved for evaluation, which gives us a held-out metric drawn from the
training distribution in addition to an external one.

\begin{table}[t]
\centering
\caption{The pinned recipe. All \nDistinctRuns{} distinct full-budget runs share
every architectural and data row (except the single 1M-batch pair and the four
Wave-12 transferred-LR runs); arms in the core comparison differ only in the
optimizer's factor choices.}
\label{tab:recipe}
\footnotesize\setlength{\tabcolsep}{6.0pt}
\begin{tabular}{llll}
\toprule
Model & & Training & \\
\midrule
Layers / width & 32 / 512 & Global batch & 524{,}288 tok. \\
MLP width & 2048 (SwiGLU) & Updates & 5{,}120 \\
Heads / head dim & 8 / 64 & Tokens per run & 2{,}684{,}354{,}560 \\
Context & 4{,}096 & World size & 8 A100-40GB \\
Vocabulary & 128{,}256 & Microbatch / accum. & 4 / 4 \\
Block matrices & 224 (134.3M) & Schedule & 1024 flat, linear decay \\
Total parameters & 265{,}585{,}152 & Precision & BF16, FP32 params \\
\addlinespace
Optimizer & & Statistics & \\
\midrule
Momentum / Nesterov & 0.95 / on & Moment EMA $\beta$ & 0.938 \\
Newton--Schulz steps & 10 & Damping $\epsilon$ & $10^{-3}$ \\
NS arithmetic & BF16 ops, FP32 out & Sampled positions & 286 per seq. \\
Decay $\lambda$ & 0.1, shaped & Samples per update & 36{,}608 \\
Inner / outer powers & $(1/2,\,1/8)$ & Dense refresh & 4 updates, async \\
Matrix base LR & 0.02262741696 & Auxiliary AdamW LR & 0.0048 \\
\bottomrule
\end{tabular}
\end{table}

\subsection{Metrics, and which one decides}

Each run records three quantities, and they are not interchangeable.

\textbf{Selection loss} is the arm-choosing metric: a token-weighted average
over the three reserved held-out sources using the mixture's own weights,
$(3.8\,\text{DCLM}+0.25\,\text{StarCoder}+0.055\,\text{Proof-Pile})/4.105$.
Each evaluation covers about 786{,}000 tokens the optimizer has never seen.

\textbf{C4 loss} is the Paloma \texttt{c4\_en} validation loss over all 236
windows, an external set held fixed across seeds. Because the reserved held-out
windows are drawn using the training seed, selection losses are comparable only
\emph{within} a seed, whereas C4 is comparable across them. The difference is
large enough to mislead: our two baselines differ by $-0.048$ nats in windowed
selection loss while differing by only $+0.001$ in windowed C4, with held-out
components moving in opposite directions. That is the signature of a different
evaluation sample, not a different model. Every delta we report is therefore
taken against its \emph{own seed's} (and, where implementation waves introduce
host-polling drift, its \emph{own wave's}) baseline, and cross-seed statements
use C4.

\textbf{Training loss} is recorded at every update --- 5{,}120 points per run
against 22 evaluations --- and is by far the densest statistic available. We use
it for resolution and consistency, never for choosing an arm. It scores data the
optimizer has just consumed, and the ranking it produces does disagree with the
held-out ranking: the dense lagged arm beats the leading arm on final-window
training loss by $0.000664$ nats with a standard error of $0.000043$, fifteen
times the error and not noise, yet it is worse on both held-out metrics. Dense
left factors fit consumed tokens slightly better while generalizing slightly
worse, which is exactly the pattern training loss cannot report.

Evaluations run at updates $1,251,501,\dots,5001,5120$: 22 per run, of which the
11 in the second half of the budget form our primary reporting window. We
report that 11-point mean rather than a single final evaluation, because the
final point alone proved too noisy to rank arms --- it was the first statistic
this project recorded and the first one it had to stop trusting.

\subsection{How small effects are resolved}
\label{sec:protocol-stats}

The effects here are a few thousandths of a nat, so the statistical question is
not incidental. Exact pairing is what makes them measurable. Within a seed,
every arm consumes identical tokens in identical order, so we difference
training loss update by update and average the differences, rather than
differencing averages. The resulting paired series has an effective standard
error of \emph{0.000038 to 0.000064} nats for arms whose difference series is
uncorrelated --- against candidate effects near 0.005, a margin of two orders
of magnitude.

Three arms are exceptions, and their exceptionality is informative rather than
inconvenient. \texttt{muon} ($0.000160$), \leftall{} ($0.000115$) and
\texttt{inner\_outer\_all\_diag} ($0.000306$) have strongly autocorrelated
paired differences, with lag-1 coefficients $0.64$, $0.57$ and $0.82$, and they
are also the only arms whose delta moves materially with the choice of window.
Where an arm's trajectory is genuinely diverging from the baseline rather than
sitting at a stable offset, no single window summarizes it, and the
autocorrelation is the signal that this is happening. We inflate the standard
error by the lag-1 factor throughout and report the five 536{,}870{,}912-token
windows separately whenever an arm's delta is not flat.

Two consequences for reading this paper. A difference of 0.0005 nats between
two candidates in a paired training window is resolvable; the same difference
between two 11-point held-out means is not, and we do not rank arms on it.
And pairing does not survive across seeds: differencing a seed-2 run against a
seed-0 run by update index gives a standard error of $0.001867$, forty times the
within-seed figure, because the series is then dominated by batch composition
rather than by the optimizer.

\subsection{What we count as a result}

The study ran in waves, and not all of them produced training. Three waves ---
one arm-geometry wave, one numerical-path wave, and one production attempt ---
stopped at correctness gates or died before their first optimizer update; their
failures are recorded in Appendix~\ref{app:runs} and contribute no optimizer
numbers. Separately, several completed waves re-ran configurations
identical to earlier ones in order to measure systems changes, and their
recorded losses are \emph{bitwise identical} to the originals across all
5{,}120 updates. Comparing full training-loss trajectories rather than trusting
wave labels, the \nRunDirs{} completed run directories hold \nDistinctRuns{}
distinct trajectories, with \nDuplicateGroups{} identical groups covering
\nDuplicateDirs{} directories (all four runs of the uncontended eight-A100
\texttt{wave\_sym\_fastlag} wave are distinct trajectories, as asynchronous host
\texttt{eigh} polling diverges after step~\symWaveExactSteps{};
Table~\ref{tab:evidence} in Section~\ref{sec:limitations}). Every count in this
paper is of distinct trajectories. Reporting deterministic repeats as
replications would be the single easiest way to overstate this study, and we
flag it here because the temptation is structural.
